\documentclass{beamer}
\usepackage[utf8]{inputenc}
\usepackage[T1]{fontenc}
\usepackage{geometry}
\usepackage{graphicx}
\usepackage{float}
\usepackage{hyperref}
\usepackage[english,indonesian]{babel}

\title{Optimasi Pemrosesan Template LaTeX Menggunakan Pendekatan Berbasis AI}
\author{Academic Support Agent}
\date{\today}

\begin{document}
\maketitle

\begin{abstract}

\end{abstract}

\chapter{PENDAHULUAN}


\chapter{TINJAUAN PUSTAKA}


% Original content (commented out):
% \documentclass[8pt,aspectratio = 1610]{beamer}

\usetheme[sectionpage=none, subsectionpage=progressbar]{metropolis}
%\usepgfplotslibrary{dateplot}
\useoutertheme[subsection=false,footline=authortitle]{miniframes}
\setbeamertemplate{footline}[frame number]%Numerando as Páginas
\hypersetup{pdfpagemode=FullScreen}%Visualizar em Tela cheia
\hypersetup{pdfpagelayout=SinglePage}%Layout da pagina
%\usetheme[
%%% Modelo de apresentação modificado para a paleta de cores da FGR
%%% option passed to the outer theme
%    progressstyle=fixedCircCnt,   % fixedCircCnt, movingCircCnt (moving is deault)
%  ]{Feather}
  
% If you want to change the colors of the various elements in the theme, edit and uncomment the following lines

% Change the bar colors:
%\setbeamercolor{Feather}{fg=white!20,bg=blue}

% Change the color of the structural elements:
%\setbeamercolor{structure}{fg=red}

% Change the frame title text color:
%\setbeamercolor{frametitle}{fg=blue}

% Change the normal text color background:
%\setbeamercolor{normal text}{fg=black,bg=gray!10}

%-------------------------------------------------------
% INCLUDE PACKAGES
%-------------------------------------------------------

\usepackage[utf8]{inputenc}
\usepackage[english,brazil]{babel}
\usepackage[T1]{fontenc}
\usepackage{helvet}
\usepackage[num]{abntex2cite}%Citações padrão ABNT}
\citebrackets[]
\usepackage{pgfplots}
\usepackage{amsmath}


\usepackage{pdflscape}
\usepackage{afterpage}
\usepackage{capt-of}
\usepackage{booktabs}
\usepackage{array}
\usepackage{multirow}
\usepackage{longtable}
\usepackage{siunitx} % for 'S' column type
\newcolumntype{T}[1]{S[table-format=#1]}
\newcommand\mc[1]{\multicolumn{1}{c}{#1}}
\usepackage{graphicx}
\usepackage{caption}

%-------------------------------------------------------
% DEFFINING AND REDEFINING COMMANDS
%-------------------------------------------------------

% colored hyperlinks
\newcommand{\chref}[2]{
  \href{#1}{{\usebeamercolor[bg]{Feather}#2}}
}

%-------------------------------------------------------
% INFORMATION IN THE TITLE PAGE
%-------------------------------------------------------
 %\institute{\textcolor{verde}{Programa de Pós-Graduação Profissional em Engenharia Elétrica (PPEE)}}
\titlegraphic{
\begin{picture}(0,0)
\put(285,15){\makebox(0,0)[rt]{\includegraphics[width=10cm]{images/BannerPPEE.png}}}
\put(400,-230){\makebox(0,0)[rt]{\includegraphics[width=5cm]{images/latitude (1).png}}}
\end{picture}
}
\title[] % [] is optional - is placed on the bottom of the sidebar on every slide
{ % is placed on the title page
      \textbf{Title
}
}

\subtitle[Nome do rodapé do slide]
{
    \textbf{Defesa de Mestrado Profissional\\}
    PPEE.MP.XX
}

\author[NAME]
{ \normalsize     
      Aluno: Name \\
      Orientador: Prof. Dr. Name \\
      Coorientador: Prof. Dr. Name
}

\institute[]
{\small
      Universidade de Brasília\\
      Faculdade de Tecnologia (FT) \\
      Departamento de Engenharia Elétrica (ENE) \\
      Programa de Pós-Graduação Profissional em Engenharia Elétrica (PPEE) \\

  
  %there must be an empty line above this line - otherwise some unwanted space is added between the university and the country (I do not know why;( )
}

\date{\today}

%-------------------------------------------------------
% THE BODY OF THE PRESENTATION
%-------------------------------------------------------

\begin{document}

%-------------------------------------------------------
% THE TITLEPAGE
%-------------------------------------------------------

% % this is the name of the PDF file for the background
\begin{frame}[plain,noframenumbering] % the plain option removes the header from the title page, noframenumbering removes the numbering of this frame only
  \titlepage % call the title page information from above
\end{frame}


\begin{frame}{Agenda}{}
\tableofcontents

\end{frame}

%-------------------------------------------------------
\section{Introdução}
%-------------------------------------------------------
\subsection{Text}
\begin{frame}{Introdução}{Motivações e Problema}
%-------------------------------------------------------
\textbf{Motivação}
        \begin{itemize}
             \item Text.
             \item Text. 
             \item Text  
         \end{itemize}
\textbf{Problema}
        \begin{itemize}
            \item Text?
            \begin{itemize}
                  \item Text. 
                 \item \textbf Text.
             \end{itemize} 
         \end{itemize} 
%    \textbf{Contribuições do Trabalho}
%        \begin{itemize}
%            \item Identificar as percepções dos gestores públicos brasileiros no uso de computação em nuvem quanto:
%            \begin{itemize}
%                 \item \textbf{Maturidade no uso} de computação em nuvem; \\
%                 \item \textbf{Domínios de recomendação para gestão de riscos} em computação em nuvem; \\
%                 \item \textbf{Riscos e Controles envolvidos} no uso de computação em nuvem; \\
%             \end{itemize} 
%         \end{itemize} 
\end{frame}


\begin{frame}{Introdução}{Objetivo}
\textbf{Objetivo}


      \begin{itemize}
      \item Geral: Text.
      \begin{itemize}
        \item Específico 1: Text;\\
        \item Específico 2: Text;\\
        \item Específico 3: Text;\\
        \item Específico 4: Text;\\
        \item Específico 5: Text;\\
      \end{itemize}
    \end{itemize}
    
\textbf{Classificação da Pesquisa}
  \begin{itemize}
             \item Objetivos: \textbf{Text};
             \item Abordagem: \textbf{\textbf{Text}};
             \item Procedimentos Técnicos: \textbf{Text};
             \item Temporalidade: \textbf{Text};
  \end{itemize}
  

         
\end{frame}

%-------------------------------------------------------
\section{Revisão da literatura}
%-------------------------------------------------------
\subsection{Text, Text e Trabalhos correlatos}
\begin{frame}{Revisão da literatura}{}
%-------------------------------------------------------
\begin{block}{}

\textbf{Text:}
    \begin{itemize}
         \item Text;
         \item Text;
         \item Text;
    \end{itemize} 
 \vspace{5pt} 
\textbf{Text:}
    \begin{itemize}
         \item Text;
         \item Text;
         \item \textbf{Text} 
         \item Text;
    \end{itemize}
     \vspace{5pt} 
 \textbf{Trabalhos correlatos:}
    \begin{itemize}
         \item Text;
    \end{itemize}

\end{block}
\end{frame}

\begin{frame}{Revisão da literatura}{}
\begin{block}{}
\textbf{Text:}
 \begin{itemize}
     \item Text;
     \item Text;
 \end{itemize}
 \include{Tabelas/Trabalhos}

\end{block}
\end{frame}

%-------------------------------------------------------
\section{Metodologia}
\subsection{Amostra e Entrevista Semiestruturada}
%\begin{frame}{Metodologia}{}
%-------------------------------------------------------

%\begin{block}{Amostra}
%    \begin{minipage}{1.0\textwidth}
% \begin{minipage}[c]{0.49\textwidth}
% \centering
% \includegraphics[width=0.85\textwidth]{images/timeline.png}
 % \rule{6.4cm}{3.6cm}
% \captionof{figure}{\textit{Timeline} do Estudo Longitudinal (Fig. 3.1)}
% \end{minipage}
% \hfill
% \begin{minipage}[c]{0.49\textwidth}
% \centering
% \include{Tabelas/amostra}
%  \end{minipage}
%\end{minipage}
%\end{block}
%\end{frame}

\begin{frame}{Metodologia}{}
     \begin{block} {Text}
     
          
          \end{block}
\end{frame}

\begin{frame}{Metodologia}{}
     \begin{block} {Text}
      \vspace{5pt} 
            Text~\cite{enisa2012}, NIST~\cite{Jansen2011} e CSA~\cite{CloudSecurityAlliance2017}, fText:
            \begin{itemize}
            \item Text;
            \item Text;
            \item Text;
            \item Text;
            \end{itemize}
     \end{block}
     \begin{block}{Análise de Conteúdo}
        \vspace{5pt} 
    Análise de Conteúdo de acordo com o Text \cite{Bardin}:
       \begin{itemize}
            \item Text;
           \item Text
           \item Text
       \end{itemize}
     \end{block}
\end{frame}



%-------------------------------------------------------
\section{Análise dos Resultados}
\subsection{Unidade de Registro / Contexto}
\begin{frame}{Análise dos Resultados}{Unidade de Registro e de Contexto}
%-------------------------------------------------------
Text;
\begin{itemize}
    \item Categoria 1: Text;
    \item Categoria 2: Text;
    \begin{itemize}
        \item Subcategoria 2.1: Text;
        \item Subcategoria 2.2: Text;
    \end{itemize}
    \item Categoria 3: Text;
    \item Categoria 4: Text:
    
\end{itemize}
\end{frame}

\begin{frame}{Análise dos Resultados}{Text}
 \textbf{Categoria 1 – Text;}
\include{Tabelas/organizacao}
\end{frame}

\begin{frame}{Análise dos Resultados}{Análise de Resultados}
%-------------------------------------------------------
\textbf{Categoria 2 – Text}
     \begin{itemize}
        \item Text:
        \begin{itemize}
            \item Text;
            \item \textbf
 Text;
        \end{itemize}
        \item Text
        \begin{itemize}
            \item \textbf Text;
        \end{itemize}
     \end{itemize}

\textbf{Categoria 3 – Text}\\
\vspace{5pt}
Text:
\begin{itemize}
    \item Text;
    \item Text.
    \item Text;
    \item Text;
    \item Text; 
    \item Text;
    \item Text.
\end{itemize}
\end{frame}


\begin{frame}{Análise dos Resultados}{Análise de Resultados}
%-------------------------------------------------------
\textbf{Categoria 4 – Text}\\
Text 

\begin{figure}[ht] % utilizar o h! ou htpb para posicionar a imagem exatamente onde foi colocada no código fonte
	\begin{center}
 
images/Figura.jpeg

    \caption{Text (Fig. x.x).}
    \end{center}
\end{figure}
\begin{itemize}
    \item Text:
    \begin{itemize}
        \item Text.
        \item Text.
    \end{itemize}
\end{itemize} 

\end{frame}

\begin{frame}{Análise dos Resultados}{Análise de Resultados}
%-------------------------------------------------------
\textbf{Categoria 4 – Text}\\
Text


\begin{itemize}
    \item Text:
    \begin{itemize}
        \item \textbf Text \cite{Da_Republica_undated-vx} .
        \item \textbf Text.
    \end{itemize}
\end{itemize} 
\end{frame}

\begin{frame}{Análise dos Resultados}{Análise de Resultados}
%-------------------------------------------------------
\textbf{Text – Text}\\
Text


\end{frame}

\begin{frame}{Análise dos Resultados}{Análise de Resultados}
%-------------------------------------------------------
\textbf{Text}\\
Text 

\begin{figure}[ht] % utilizar o h! ou htpb para posicionar a imagem exatamente onde foi colocada no código fonte
	\begin{center}

    \includegraphics[width=0.80\textwidth]{images/controles.png}\\
    \caption{Text (Fig. x.x).}
    \end{center}
\end{figure}

\begin{itemize}
    \item Text:
    \begin{itemize}
        \item Text.
        \item Text.
        \item Text.
        \item Text.
    \end{itemize}
\end{itemize}

\end{frame}



%-------------------------------------------------------
\section{Conclusão}
\subsection{Conclusão, Trabalhos Futuros e Publicações}
\begin{frame}{Conclusão}{Conclusão}
%-------------------------------------------------------

  \begin{block}{Conclusão}
\begin{itemize}
    \item Text.
    \vspace{5pt}
    \item Text.
    \vspace{5pt}
    \item Text
    \vspace{5pt}
    \begin{itemize}
        \item Text. 
        \vspace{5pt}
        \item Text.
    \end{itemize}
\end{itemize}
  \end{block}
\end{frame}

\begin{frame}{Conclusão}{Trabalhos Futuros}
%-------------------------------------------------------

\begin{block}{Trabalhos Futuro}
\vspace{5pt}

    \begin{itemize}
        \item Text.
        \vspace{5pt}
        \item Text.
        \vspace{5pt}
        \item Text.
    \end{itemize}
  \end{block}
\end{frame}

\begin{frame}{Conclusão}{Publicações}

  \begin{block}{Publicações}
    
\begin{center}
\fontsize{7}{7}\selectfont
\begin{longtable}{p{2.0cm}p{10.5cm}}
\multicolumn{2}{c}{} 
\endfirsthead
\textbf{Título 1:} &  \textit{\textbf{Text}}~\cite{Alves2021}.\\
Autores: Text\\
RIS: & Text. \\
Publicação: &  Text.\\
url: & \url{Text}\\
DOI: &  \url{Text} \\\\\\

\textbf{Título 2:} &  \textit{\textbf{Text}}~\cite{Tabosa2021}.\\
Autores: &  Text\\
Conferência: &  Text.\\
Publicação: & Text \\
url: & \url{Text}\\
DOI: &  \url{Text}
\end{longtable}
\end{center}

  \end{block}
\end{frame}



%================================================
%---------------Elementos Pós-textuais
%================================================
{\setbeamercolor{palette primary}{fg=black, bg=white}
\begin{frame}[standout]
Obrigado! \\
Perguntas e Comentários
\end{frame}
} 


%================================================
%----------------BIBLIOGRAFIA
%================================================
\section*{Bibliografia}
\begin{frame}[allowframebreaks]{Referências bibliográficas}
%\nocite{*}

\setbeamertemplate{bibliography item}[text] 


\bibliography{bibliografia}
% \printbibliography

\end{frame}
%================================================
%----------------FIM DO TEMPLATE
%================================================

%\begin{frame}[plain,noframenumbering]
%  \finalpage{Thank you for using Feather Beamer Theme!}
%\end{frame}

\end{document}

\end{document}